\documentclass[../main.tex]{subfiles}

\begin{document}

\onlyinsubfile{
    \renewcommand{\contentsname}{ОГЛАВЛЕНИЕ}
    \setcounter{tocdepth}{3}
    \tableofcontents
}

\refstepcounter{chapter}
\chapter*{\MakeUppercase{Глава \arabic{chapter}\\ Оценка, ограничения и перспективы развития генеративных моделей}}\label{chapter:evaluation_and_perspectives}
\addcontentsline{toc}{chapter}{Глава \arabic{chapter} Оценка, ограничения и перспективы развития генеративных моделей}

Анализ ключевых архитектур, представленный в предыдущей главе, демонстрирует значительный прогресс в области генеративного дизайна молекул. Однако само по себе создание новых моделей не является достаточным условием для научного развития. Не менее важными являются стандартизация их оценки, глубокое понимание фундаментальных ограничений существующих подходов и четкое видение перспективных направлений исследований. Данная глава посвящена критическому анализу этих трех аспектов, синтезируя выводы из передовых научных работ и формируя целостную картину текущего состояния и будущего генеративной медицинской химии.

% Содержимое подраздела 3.1 - Стандартизация оценки: бенчмарки и метрики


\section{Стандартизация оценки: бенчмарки и метрики}
\label{sec:standardization}

Одной из центральных проблем, тормозивших развитие области генеративных моделей, было отсутствие единого стандарта для их сравнения. Исследователи использовали различные наборы данных, методы их предобработки и метрики качества, что делало невозможным объективное сравнение эффективности новых архитектур по принципу <<яблоки с яблоками>>. Это приводило к несопоставимым, зачастую невоспроизводимым результатам и затрудняло оценку реального прогресса \cite{polykovskiy2020molecular}.

Ответом на этот вызов стала разработка комплексной платформы с открытым исходным кодом MOSES (Molecular Sets) \cite{polykovskiy2020molecular}. Данная платформа не просто предложила новый метод, а создала инфраструктуру для стандартизированного обучения и оценки, состоящую из трех ключевых компонентов:
\begin{enumerate}
    \item \textbf{Стандартизированный набор данных.} Авторы MOSES тщательно отфильтровали коллекцию ZINC Clean Leads, создав датасет с четко определенным разделением на обучающую, тестовую и, что особенно важно, \textit{scaffold test} выборки. Последняя содержит молекулы с структурными каркасами, отсутствующими в обучающей выборке, что позволяет оценить способность модели к генерализации и созданию принципиально новых химических структур.
    \item \textbf{Реализация базовых моделей.} Платформа включает эталонные реализации нескольких ключевых архитектур (CharRNN, VAE, AAE, JTN-VAE), которые служат <<точкой отсчета>> (reference point) для всех будущих моделей.
    \item \textbf{Комплексный набор метрик.} MOSES формализовал и популяризовал набор метрик для всесторонней оценки генеративных моделей, который впоследствии стал стандартом де-факто в большинстве последующих работ \cite{mouchlis2021advances, kim2021comprehensive}.
\end{enumerate}

Ключевые метрики, предложенные в MOSES, можно разделить на несколько категорий. Базовые метрики качества генерации включают:
\begin{itemize}
    \item \textbf{Validity (Валидность):} Процент сгенерированных строк (например, SMILES), которые являются синтаксически и семантически корректными и могут быть преобразованы в валидную химическую структуру.
    \item \textbf{Uniqueness (Уникальность):} Процент уникальных молекул среди всех валидных сгенерированных. Низкое значение этого показателя может свидетельствовать о коллапсе мод (mode collapse), когда модель генерирует лишь небольшое подмножество известных ей структур.
    \item \textbf{Novelty (Новизна):} Процент уникальных валидных молекул, отсутствующих в обучающем наборе данных. Эта метрика напрямую измеряет способность модели создавать новые, ранее не виденные соединения.
    \item \textbf{Internal Diversity (IntDiv, Внутреннее разнообразие):} Среднее попарное несходство (обычно 1 - Tanimoto similarity) молекул внутри сгенерированного набора. Высокое значение указывает на то, что модель генерирует химически разнообразные структуры.
\end{itemize}

Однако наиболее важным вкладом MOSES стала популяризация метрик, оценивающих сходство распределений. Среди них центральное место занимает \textbf{Fréchet ChemNet Distance (FCD) \cite{polykovskiy2020molecular}. FCD измеряет расстояние Фреше между распределениями активаций предпоследнего слоя нейросети ChemNet (предобученной на предсказание биологической активности) для сгенерированного и референсного наборов молекул. Низкое значение FCD указывает на то, что сгенерированные молекулы не только структурно, но и с точки зрения потенциальных биологических свойств похожи на реальные данные. Эта метрика стала мощным интегральным показателем, позволяющим оценить, насколько хорошо генеративная модель уловила сложное многомерное распределение химически осмысленных структур.

% Содержимое подраздела 3.2 - Фундаментальные ограничения строковых представлений

\section{Фундаментальные ограничения строковых представлений}
\label{sec:string_limitations}

Несмотря на популярность и удобство, строковые представления молекул, в первую очередь SMILES, накладывают серьезные фундаментальные ограничения на генеративные модели. Эти ограничения можно разделить на две категории: синтаксические и семантические.

Первая проблема -- \textbf{синтаксическая корректность}. Модели, работающие на символьном уровне, такие как RNN, генерируют SMILES-строку символ за символом, пытаясь выучить сложную <<грамматику>> этого языка. Однако этот процесс часто приводит к генерации синтаксически некорректных строк, которые не соответствуют ни одной реальной молекуле. Эта проблема привела к разработке подходов, которые работают не со строками, а напрямую с графовой структурой молекул. Ярким примером является \textbf{Junction Tree Variational Autoencoder (JT-VAE)} \cite{jin2018junction}. Эта модель решает проблему валидности <<по построению>>. Вместо генерации атомов, JT-VAE оперирует химически осмысленными строительными блоками (кольцами, простыми связями). Генерация происходит в два этапа: сначала создается древовидный каркас из этих блоков, а затем они собираются в итоговый молекулярный граф. Поскольку на каждом шаге используются только валидные химические компоненты, все сгенерированные молекулы гарантированно являются 100\% химически корректными.

Вторая, более глубокая проблема -- \textbf{семантическая неадекватность}. SMILES представляет молекулу на уровне <<букв>> (атомов), в то время как химики мыслят на уровне <<слов>> (функциональных групп, фрагментов). Кроме того, SMILES имеет серьезные проблемы с однозначным и инвариантным представлением хиральности -- ключевого свойства для биологической активности. Небольшие изменения в порядке обхода графа могут кардинально изменить SMILES-строку и представление стереоцентров. Эту проблему решает нотация \textbf{fragSMILES} \cite{mastrolorito2025fragsmiles}. Она представляет собой принципиальный переход от атомно-уровневого представления к фрагментному. Молекула сначала декомпозируется на химически осмысленные фрагменты, которые и становятся токенами (<<словами>>) в новой строке. Информация о хиральности кодируется явно и привязывается к точкам соединения фрагментов, что делает ее инвариантной к порядку обхода. Эксперименты показывают, что модели, обученные на fragSMILES, генерируют значительно более качественные и химически релевантные структуры, особенно в задачах, требующих точного контроля над стереохимией.

% Содержимое подраздела 3.3 - Перспективы развития: за пределами RNN и SMILES


\section{Перспективы развития: за пределами RNN и SMILES}
\label{sec:future_perspectives}

Анализ ограничений классических подходов позволяет очертить три магистральных направления, по которым движется современная наука в области генеративного дизайна молекул: переход к более мощным архитектурам, переход к более естественным представлениям данных и, наконец, переход от 2D-структур к 3D-геометрии.

\textbf{Тренд №1: От RNN к Трансформерам и большим языковым моделям.}
Рекуррентные нейронные сети, долгое время доминировавшие в задачах обработки последовательностей, постепенно уступают место архитектуре Трансформер. Благодаря механизму внимания (attention), трансформеры способны улавливать дальние зависимости в последовательностях, что делает их более мощным инструментом для моделирования сложных SMILES-строк. Эта тенденция нашла отражение в современных фреймворках, таких как \textbf{REINVENT 4}, где трансформерная модель была добавлена как опция для задач локальной оптимизации молекул \cite{loeffler2024reinvent}. Последующие работы показали, что обучение с подкреплением (RL) особенно эффективно именно в связке с трансформерами для целенаправленной генерации \cite{he2024evaluation}. Апогеем этого тренда стало появление <<фундаментальных моделей>> (foundation models), таких как \textbf{GP-MOLFORMER}, обученных на миллиардах молекул. Эти модели, заимствуя парадигму из обработки естественного языка (NLP), демонстрируют беспрецедентную универсальность и производительность \cite{ross2025gp}.

\textbf{Тренд №2: От строк к графам.}
Как было показано ранее, строковые представления являются неестественными для молекул, которые по своей природе являются графами. Поэтому вторым ключевым трендом является все более широкое применение графовых нейронных сетей (GNN). Модель \textbf{JT-VAE} была одним из пионеров этого направления. Обширные обзорные работы подтверждают, что GCN (Graph Convolutional Networks) и другие варианты GNN сегодня являются передовой технологией для широкого спектра задач, от предсказания свойств до de novo генерации, так как они способны извлекать признаки напрямую из топологии молекулы, минуя <<бутылочное горлышко>> строкового представления \cite{sun2020graph}.

\textbf{Тренд №3: От 2D к 3D.}
Наиболее передовым и сложным направлением является переход от генерации 2D-графов к генерации полноценных 3D-конформаций молекул. 3D-структура напрямую определяет биологическую активность, и ее моделирование является Святым Граалем вычислительной химии. На этом фронте доминирующим подходом стали \textbf{диффузионные модели}. Такие модели, как \textbf{GeoDiff}, обучаются обращать процесс постепенного зашумления 3D-координат атомов, генерируя физически правдоподобные и геометрически корректные структуры с соблюдением фундаментальных симметрий \cite{xu2022geodiff}. Как и в случае с 2D-моделями, эта бурно развивающаяся область уже столкнулась с необходимостью стандартизации, что привело к появлению комплексных бенчмарков для сравнения 3D-генераторов \cite{qin2025comprehensive}. Финальным аккордом этого тренда можно считать появление моделей вроде \textbf{BindGPT}, которые объединяют в себе все три направления: используют архитектуру трансформера (LLM), работают с единым представлением, объединяющим 2D-граф (SMILES) и 3D-координаты, и интегрируют обучение с подкреплением для целенаправленной генерации молекул с высокой аффинностью к белковой мишени. Это демонстрирует путь к созданию единых, универсальных и мощных инструментов для полностью автоматизированного дизайна лекарств \cite{zholus2025bindgpt}.


\onlyinsubfile{
    \printbibliography[heading=bibintoc, title={Библиографический список}]
}

\end{document}
